\section{Related Work}
\label{sec:related_work}

\para{Believable simulation is not the same as distributional validity.} Generative agent work shows that LLM-based systems can produce coherent, believable behavior in interactive environments \cite{park2023generative}. That line of work is important because it demonstrates process-level plausibility and coordination. However, believable behavior in a sandbox does not establish that the model captures the empirical distribution of human outcomes. Our focus is narrower: we ask whether conditioning can widen behavior while preserving fidelity to known regularities.

\para{Human-study replication motivates our task design.} Aher et al.\ introduce Turing Experiments as a way to test whether language models can reproduce established findings from behavioral studies \cite{aher2023using}. Xie et al.\ show that GPT-4-class agents can align with some human trust-game regularities when equipped with careful prompting and agent structure \cite{xie2024trust}. We build on this evaluation tradition, but we focus on two distinct levels of evidence: benchmarked social judgment and population-style policy variation. Unlike prior work that treats successful replication on one task as broad evidence of simulation, we explicitly separate local fidelity from coverage.

\para{Population alignment remains a central challenge.} Santurkar et al.\ demonstrate that demographic steering does not reliably recover the opinion distributions of real demographic groups \cite{santurkar2023opinions}. Lin argues that LLMs should be treated as pragmatic simulation tools rather than replacements for human participants, and warns against reading alignment as explanation \cite{lin2025fallacies}. Our work operationalizes this conceptual caution. We ask not only whether the model looks plausible, but whether latent conditioning changes the response distribution in a way that supports stronger simulation claims.

\para{Recent evidence points to collapse toward an average, rationalized person.} Several studies now report systematic distortions in LLM-based human simulation. Sorokovikova et al.\ find that models can express personality-like variation, but the realism and stability of these traits remain open \cite{sorokovikova2024bigfive}. Liu et al.\ show that strong LLMs assume people are more rational than they really are \cite{liu2025rational}. Chen et al.\ report that long-horizon simulation benchmarks reveal hyper-activity, persona homogenization, and utopian bias \cite{chen2026omnibehavior}. Our results fit this picture: conditioning increases variety, but the variation remains constrained inside a narrow and overly cooperative family of responses.
